% Section 7: Findings — Where Formalization Breaks Down

\section{Findings: Where Formalization Breaks Down}
\label{sec:findings}

\subsection{Category 1: Open-Textured Concepts}
\label{sec:open-texture}

Certain legal terms resist formalization:

\begin{table*}[t]
\centering
\begin{tabular}{@{}lll@{}}
\toprule
Term & Source & Why It Resists \\
\midrule
``Defined investment policy'' & AIFMD Art. 4 & What counts as ``defined''? \\
``Expertise, experience and knowledge'' & MiFID II Annex II & Inherently qualitative \\
``Adequate'' assessment & MiFID II & Procedural adequacy is judgment \\
``Commercially reasonable'' & LPA boilerplate & Context-dependent standard \\
``Material adverse change'' & Side letters & Requires interpretation \\
``Best efforts'' / ``Reasonable efforts'' & Various & Gradations without bright lines \\
\bottomrule
\end{tabular}
\caption{Open-textured concepts resisting formalization}
\label{tab:open-texture}
\end{table*}

These concepts cannot be encoded without making interpretive choices that may not match regulatory or judicial interpretation. Hart's concept of ``open texture'' \citep{hart1961concept} suggests this is inherent to law, not merely a drafting deficiency.

\subsection{Category 2: Jurisdictional Variation}
\label{sec:jurisdiction}

The same directive is transposed differently across jurisdictions:

\begin{itemize}
  \item Luxembourg RAIF Law vs.\ Irish QIAIF regime
  \item CSSF guidance vs.\ FCA guidance on same AIFMD provision
  \item US LP investors subject to Volcker Rule overlay
  \item Cayman feeder funds subject to different AML requirements
\end{itemize}

A specification for ``Luxembourg fund structures'' cannot capture these variations without becoming a multi-jurisdictional model of prohibitive complexity. Even within the EU, national gold-plating creates divergence.

\subsection{Category 3: Temporal Dynamics}
\label{sec:temporal}

Fund structures evolve:

\begin{itemize}
  \item Amendment 17 may contradict Amendment 3
  \item An investor who was ``professional'' at subscription may no longer qualify
  \item Regulatory guidance changes; old structures may be ``grandfathered''
  \item Investment restrictions may change between closings
\end{itemize}

A specification is a snapshot. The fund is a process. Verification would need to occur at each modification, with version control of the specification itself. This is tractable but multiplies complexity.

\subsection{Category 4: The Specification-Reality Gap}
\label{sec:gap}

Even a ``complete'' specification verifies only:

\begin{itemize}
  \item That the \textit{encoded} rules are consistent
  \item That \textit{inputs} (if accurate) produce correct \textit{outputs}
\end{itemize}

It does NOT verify:

\begin{itemize}
  \item That the encoding matches legal meaning (the interpretation problem)
  \item That inputs are accurate (the data quality problem)
  \item That the law will be interpreted as encoded (the adjudication problem)
\end{itemize}

Formal verification is verification \textit{relative to a model}. The model is not the law. A fund could pass all formal checks and still face regulatory enforcement if the specification misunderstood the legal requirement.

\subsection{Category 5: The Incompleteness of Written Law}
\label{sec:incompleteness}

Much of fund practice operates in a space beyond written law:

\begin{itemize}
  \item Market convention (what a ``market standard'' side letter contains)
  \item Regulatory practice (what CSSF actually accepts vs.\ what law says)
  \item Professional judgment (what a ``reasonable'' board decision looks like)
\end{itemize}

These unwritten norms cannot be specified because they are not specified in the source material. They exist in practice, transmitted through precedent and professional socialization \citep{macaulay1963non}.
